\documentclass[12pt,a4paper]{article}
\usepackage[UTF8]{ctex}
\usepackage{graphicx}
\usepackage{geometry}
\usepackage{ulem}
\usepackage{enumitem}
\usepackage{titlesec}
\usepackage{multicol}
\usepackage[export]{adjustbox}
\usepackage{amsmath}  % 推荐用于更好的数学排版
\usepackage[utf8]{inputenc}
\usepackage{CJKutf8}
\usepackage{fancyhdr}
\usepackage{lastpage}
\usepackage{ulem}
\usepackage{tabularx}
\usepackage{float}
\usepackage{siunitx}
\usepackage{multirow} 

% 页眉页脚设置
\pagestyle{fancy}
\fancyhf{} % 清除所有页眉页脚
% 设置页眉内容
%-==========1. 填写从第二页开始 每页第一行的信息==========-
\fancyhead[C]{\small
    \makebox[4.5em][c]{实验名称：} \uline{\makebox[17em][c]{基本实验测试}} 
    \makebox[2.5em][c]{姓名：} \uline{\makebox[4em][c]{}} 
    \makebox[2.5em][c]{学号：} \uline{\makebox[5em][c]{}}
}%--------------------------------------------------------
\fancyfoot[C]{\small 第 \thepage\ 页，共 \pageref{LastPage} 页}


% 取消页眉和页脚的下划线
\renewcommand{\headrulewidth}{0pt} % 页眉横线粗细设为0
\renewcommand{\footrulewidth}{0pt} % 页脚横线粗细设为0

% 调整页眉高度
\setlength{\headheight}{50pt}

% 设置plain样式（用于第一页）
\fancypagestyle{plain}{
    \fancyhf{} % 清除所有页眉页脚
    \renewcommand{\headrulewidth}{0pt} % 隐藏页眉横线
    \fancyfoot[C]{第 \thepage\ 页，共 \pageref{LastPage} 页} % 保留页脚
}
% 设置页面边距
\geometry{left=2.54cm,right=1.5cm,top=2.54cm,bottom=2.54cm}

% 设置正文字体大小为小四
\renewcommand{\normalsize}{\fontsize{12pt}{16pt}\selectfont}
% 设置章节格式
\titleformat{\section}{\fontsize{16pt}{22pt}\selectfont\bfseries}{\thesection}{1em}{} % 一级标题：三号
\titleformat{\subsection}{\fontsize{15pt}{20pt}\selectfont\bfseries}{\thesubsection}{1em}{} % 二级标题：小三

\begin{document}
\thispagestyle{plain}%第一页使用 plain 样式，而不是默认的 fancy 样式，第一页不再显示页眉
% 标题部分 - 两栏布局
\begin{multicols}{2}
    % 左栏：图片和实验报告标题
    % 左栏：使用空行创建前三行
    \noindent
    ~\\ % 第一行
    ~\\ % 第二行
    ~\\ % 第三行
    \hphantom{占位}% 使用空占位符
    \raisebox{-0.5\baselineskip}% 图片向下移动半行
    {
\includegraphics[width=1.85in,height=0.58in]{ZJU.jpg}}
    \hspace{.5em}% 图片和标题之间的间距
    \raisebox{0pt}[\height][0pt]% 标题底部对齐
    {\textbf{\fontsize{26pt}{30pt}\selectfont 实验报告}}
    
    \columnbreak % 换到右栏
    
    % 右栏：个人信息
    % 右栏：个人信息 - 靠右对齐 
%-==========2. 填写个人信息==========-
    \begin{flushright}
        \small
        \makebox[3em][l]{专业：} \uline{\makebox[9em][c]{法语-电子科学与技术}} \\
        \makebox[3em][l]{姓名：} \uline{\makebox[9em][c]{}} \\
        \makebox[3em][l]{学号：} \uline{\makebox[9em][c]{}} \\
        \makebox[3em][l]{日期：} \uline{\makebox[9em][c]{2026年4月13日}} \\
        \makebox[3em][l]{地点：} \uline{\makebox[9em][c]{}} \\
    \end{flushright}   
\end{multicols}
% 课程名称等信息（单栏）
%-==========3. 填写课程信息==========-

{\small
\noindent
课程名称： \uline{电子电路设计实验II} \hfill 指导老师： \uline{施红军、叶险峰、邓靖靖} \hfill 成绩： \uline{\hspace{3cm}} \\
实验名称：\uline{基本实验测试}\hfill 实验类型：\uline{设计实验}\hfill 同组学生姓名：\uline{}
}
%-==========4. 正文开始==========-
\section*{一、实验目的}

通过基本实验设计，进一步了解Arduino的使用方法和编程技巧。

\section*{二、实验任务与要求}

至少选做本章三节视频中的实验。

报告应包括实验原理、程序代码、实验结果展示和分析。




\section*{三、实验步骤、调试过程、实验数据记录}

\subsection*{3.1 呼吸灯}

编写以下代码：
\begin{verbatim}
void setup() {
  pinMode(10, OUTPUT);
}
void loop() {
  for (int i = 0; i <= 255; i++)
  {
    analogWrite(10, i);
    delay(10);
  }
  delay(50);
  for (int i = 255; i >= 0; i--)
  {
    analogWrite(10, i);
    delay(10);
  }
}
\end{verbatim}

上传程序，可以观察到led灯在“呼吸”。

\begin {figure}[H]
    \centering
    \includegraphics[width=0.6\textwidth]{1.jpg}


    \end{figure}

\begin {figure}[H]
    \centering
    \includegraphics[width=0.6\textwidth]{2.jpg}

    \end{figure}

\subsection*{3.2 串口控制}

编写以下代码：
\begin{verbatim}
int LedPin = 13;
char ch;

void setup() {
    pinMode(LedPin, OUTPUT);
    Serial.begin(9600);
    Serial.println("Program Start...");
    delay(1000);
}

void loop() {
    if (Serial.available() > 0) {
        ch = Serial.read();
        if (ch == 'A') {
            digitalWrite(LedPin, HIGH);
            Serial.println(ch);
        } else if (ch == 'B') {
            digitalWrite(LedPin, LOW);
            Serial.println(ch);
        } else {
            Serial.println("Error");
        }
    }
}
\end{verbatim}

上传程序，可以观察到：串口先初始化输出"Program Start..."串口输入A时，板载LED灯亮，串口输入B时，板载LED灯灭。
\begin {figure}[H]

    \centering
    \includegraphics[width=0.6\textwidth]{3.png}

    \end{figure}
                
\begin {figure}[H]
    \centering
    \includegraphics[width=0.6\textwidth]{4.jpg}

    \end{figure}

\begin {figure}[H]
    \centering
    \includegraphics[width=0.6\textwidth]{5.jpg}

    \end{figure}


\subsection*{3.3 超声波传感器}

编写以下代码：
\begin{verbatim}
int TrigPin = 2;
int EchoPin = 3;
float cm;

void setup() {
  Serial.begin(9600);
  pinMode(TrigPin,OUTPUT);
  pinMode(TrigPin,INPUT);
}

void loop() {
  digitalWrite(TrigPin,LOW);
  delayMicroseconds(2);
  digitalWrite(TrigPin,HIGH);
  delayMicroseconds(10);
  digitalWrite(TrigPin,LOW);

  cm = pulseIn(EchoPin,HIGH)/58.8;
  cm = (int(cm*100.0))/100.0;

  Serial.print(cm);
  Serial.print("cm");
  Serial.println();
  delay(1000);
}

\end{verbatim}

上传程序，可以在串口监视器中看到每隔1s测量的距离。


\begin {figure}[H]
\centering
    \includegraphics[width=0.6\textwidth]{6.png}

    \end{figure}




\end{document}